% STRUTTURA DEL DOCUMENTO INTERO
\begin{document}
    \frontmatter
    \begin{titlepage}
    \begin{center}
        \begin{large}
            \textbf{\myUni}\\
        \end{large}

        \vspace{10pt}

        \begin{large}
            \textsc{\myDepartment}\\
        \end{large}

        \vspace{10pt}

        \begin{small}
            \textsc{\myFaculty}\\
        \end{small}

        \vspace{40pt}
        \vfill

        \begin{figure}[htbp]
            \centering
            \includegraphics[width=0.5\textwidth]{unipd-logo} %va bene così perchè le impostazioni di compilazione prevedono che le immagini siano nella cartella "images"
        \end{figure}
        
        \vspace{25pt}
        \vfill

        \begin{large}
            \textbf{\myTitle}\\
        \end{large}

        \vspace{10pt}

        \begin{normalsize}
            \textsl{\myDegree}\\
        \end{normalsize}

        %\vspace{30pt}
        \vfill

        \begin{normalsize}
            \begin{flushleft}
                \textit{Relatore}\\
                \vspace{5pt}
                \profTitle\ \myProf
            \end{flushleft}

            % You can tweak the spacing to have professor and student names on the same line
            % useful if the page is broken by a long thesis title and you need more space
            % \vspace{-52pt}

            \begin{flushright}
                \textit{Laureando}\\
                \vspace{5pt}
                \myName \\
                \vspace{5pt}
                \textit{Matricola} \myID
            \end{flushright}
        \end{normalsize}

        %\vspace{40pt}
        \vfill

        \line(1, 0){338} \\
        \begin{small}
            \textsc{Anno Accademico \myAA}
        \end{small}
    \end{center}
\end{titlepage} % pagina iniziale con titolo e logo
    %\clearpage
%\phantomsection
%\thispagestyle{empty}

%\hfill
%\vfill

%\noindent\myName: \textit{\myTitle,}
%\myDegree,
%\textcopyright\ \myTime.
 % riga del copywright
    \input{preface/dedication} % dedica
    
\cleardoublepage
\pdfbookmark{\contentsname}{tableofcontents}
\setcounter{tocdepth}{2}
\tableofcontents

\clearpage

% Elenco delle figure
\phantomsection
\pdfbookmark{\listfigurename}{lof}
\listoffigures

\clearpage

% Elenco delle tabelle
\phantomsection
\pdfbookmark{\listtablename}{lot}
\listoftables

\cleardoublepage % indice
    \cleardoublepage

    % contenuto principale: capitoli
    \mainmatter
    % capitolo 1: introduzione
\chapter{Introduzione}
\label{cap:introduzione}

\section{Contesto macroeconomico}
In ambito macroeconomico, una \textbf{guerra commerciale} è una situazione di conflitto economico tra due o più Paesi, caratterizzata dall'adozione di misure protezionistiche, 
quali l'imposizione di dazi e altre restrizioni agli scambi, in risposta alle politiche commerciali adottate dalle controparti. Tali misure possono innescare meccanismi 
di ritorsione reciproca, alterando i flussi commerciali internazionali e alimentando l'incertezza economica. 
L'imprevedibilità di tali misure genera shock di incertezza che si ripercuotono sui mercati finanziari, contribuendo ad accrescere la volatilità dei rendimenti azionari.

In questo contesto, la guerra commerciale tra Stati Uniti e Cina nel biennio 2025-2026 ha rappresentato un caso emblematico di \emph{escalation} protezionistica, caratterizzato 
da una rapida successione di misure tariffarie e contromisure.

L'\emph{escalation} ha avuto inizio con l'avvio del secondo mandato del Presidente statunitense Donald Trump, quando l'amministrazione americana ha introdotto dazi mirati 
nei confronti di Canada, Messico e Cina, innescando una progressiva intensificazione delle tensioni commerciali e una crescente instabilità sui mercati internazionali.

Le tensioni hanno raggiunto una prima fase di massima intensità nella primavera del 2025, quando, nel mese di aprile, l'aliquota statunitense sulle importazioni cinesi ha 
raggiunto il 125\%. All'\emph{escalation} si sono anche alternati episodi di distensione, come gli accordi di maggio e ottobre 2025.

La vicenda si è successivamente estesa al piano giuridico e istituzionale, culminando nel febbraio 2026 con la pronuncia della Corte Suprema degli Stati Uniti in merito 
all'utilizzo dell'\emph{International Emergency Economic Powers Act} (IEEPA).

Un elemento caratterizzante del periodo è stata la modalità di comunicazione delle decisioni politiche: molte delle decisioni più significative non sono state annunciate 
tramite canali diplomatici o conferenze programmate, bensì attraverso ordini esecutivi inaspettati, annunci sulle piattaforme social e dichiarazioni pubbliche estemporanee, 
spesso rilasciate durante gli orari di chiusura delle principali borse mondiali.

Quando un annuncio viene diffuso a mercati chiusi o durante il fine settimana, il suo effetto non è osservabile subito: i prezzi si aggiustano alla prima seduta utile 
e le diverse \gls{asset-class} possono incorporare la notizia in momenti diversi, a seconda dei fusi orari e degli orari di negoziazione. La data dell'annuncio non coincide quindi 
necessariamente con quella in cui se ne osserva l'effetto e ciò rende necessario selezionare correttamente il momento in cui l'evento entra nell'analisi.

L'analisi assume che gli annunci tariffari siano eventi esogeni decisi indipendentemente dall'andamento dei prezzi, così che le variazioni di volatilità osservate possano 
essere attribuite all'annuncio. Un modello \gls{garch} standard non distingue le variazioni di volatilità dovute a eventi esterni dalle fluttuazioni ordinarie e non permette 
quindi di misurare l'effetto degli annunci.

\section{Obiettivo e metodologie}
Alla luce di queste considerazioni, l'elaborato si pone l'obiettivo di capire se gli annunci tariffari abbiano modificato la volatilità condizionata degli strumenti 
considerati, con quale intensità e persistenza e con quali differenze tra asset class. 
L'analisi considera gli eventi che riguardano i rapporti bilaterali tra Stati Uniti e Cina. A tal fine, l'analisi empirica considera sei strumenti esposti alle tensioni 
commerciali attraverso canali di trasmissione diversi: un indice di mercato (\textbf{S\&P 500}), titoli azionari appartenenti a settori chiave (\textbf{Apple}, \textbf{Nucor} e \textbf{General Motors}), 
l'\textbf{oro} e il mercato valutario, rappresentato dal tasso di cambio \gls{offshore} \textbf{USD/CNH}, caratterizzato da una maggiore flessibilità rispetto al tasso di cambio domestico.

Dal punto di vista metodologico, l'analisi utilizza modelli della famiglia GARCH con componente asimmetrica, nello specifico \gls{gjr-garch} ed \gls{egarch}, con errori distribuiti 
secondo una t di Student. Per stimare l'effetto degli annunci politici, vengono introdotte nell'equazione della varianza le variabili \emph{dummy}, definite su una finestra 
temporale fissa di due sedute per ciascun evento analizzato, con un criterio uniforme per tutti gli eventi.


\section{Organizzazione del testo}
La presente tesi è articolata nei seguenti capitoli: 
\paragraph{Capitolo 1}
Introduce al contesto macroeconomico e agli obiettivi dell'analisi, con una descrizione dei principali strumenti finanziari considerati e delle metodologie utilizzate.

\paragraph{Capitolo 2}
Introduce ai concetti teorici alla base dell'analisi delle serie storiche finanziarie, illustrando le principali metriche di rendimento e i cosiddetti «fatti stilizzati» della volatilità, con particolare attenzione al \emph{volatility clustering} e all'effetto asimmetrico (\emph{leverage effect}). 

\paragraph{Capitolo 3}
Presenta il quadro econometrico di riferimento: i modelli \gls{arch} e GARCH, le loro estensioni asimmetriche (GJR-GARCH ed EGARCH) e la \emph{News Impact Curve}, utilizzata per confrontare graficamente la risposta della volatilità a shock di segno e dimensione diversi.

\paragraph{Capitolo 4}
Descrive dati e metodologia: gli strumenti finanziari selezionati, la cronologia degli annunci tariffari del biennio 2025-2026, la specificazione dei modelli con le variabili \emph{dummy}, le statistiche descrittive e l'analisi esplorativa delle serie. 

\paragraph{Capitolo 5}
Presenta e discute i risultati delle stime, valutando l'impatto degli eventi esogeni sulla volatilità condizionata e la persistenza dei loro effetti. Inoltre, attraverso un'analisi delle correlazioni \emph{rolling}, esamina i co-movimenti tra le diverse asset class nei periodi di maggiore tensione commerciale.

\paragraph{Capitolo 6}
Conclude il lavoro con una sintesi dei risultati ottenuti, con un'analisi critica dei limiti della ricerca e delle possibili estensioni future.


Riguardo la stesura del testo, relativamente al documento sono state adottate le seguenti convenzioni tipografiche:
\begin{itemize}
	\item gli acronimi, le abbreviazioni e i termini ambigui o di uso non comune menzionati vengono definiti nel glossario, situato alla fine del presente documento;
	\item per la prima occorrenza dei termini riportati nel glossario viene utilizzata la seguente nomenclatura: \emph{parola}\glsfirstoccur;
	\item i termini in lingua straniera o facenti parti del gergo tecnico sono evidenziati con il carattere \emph{corsivo}.
\end{itemize} %1
    % capitolo 2 (dopo introduzione): Analisi delle serie storiche finanziarie
\chapter{Analisi delle serie storiche finanziarie}
\label{cap:analisi-serie}

Per comprendere le dinamiche dei mercati è necessario condurre un'analisi delle serie storiche, il cui scopo è tradurre le fluttuazioni dei 
prezzi in metriche misurabili di rendimento e volatilità. Questo capitolo introduce tali concetti, indispensabili per poter valutare l'impatto dei dazi.

Il capitolo analizza le modalità di calcolo dei rendimenti e definisce il concetto di volatilità, per poi concentrarsi sui principali 
comportamenti tipici (i cosiddetti «fatti stilizzati») che caratterizzano i mercati finanziari.

Lo studio di queste dinamiche serve a mostrare che i modelli lineari classici non sono sufficienti per descrivere la realtà dei mercati, 
rendendo necessario l'utilizzo dei modelli econometrici della famiglia GARCH, che vengono affrontati nei capitoli successivi.


\section{I rendimenti finanziari}
Nei mercati finanziari, il prezzo indica la quotazione puntuale di un asset, mentre il rendimento ne misura la variazione relativa nel tempo\cite{libro:investments}. 
Nell'analisi delle serie storiche, l'oggetto di studio primario non è rappresentato dai prezzi degli asset, bensì dai loro rendimenti. 
Questa scelta metodologica è motivata per due ragioni\cite{campbell:econometrics}.

La prima ragione è che il rendimento fornisce un riassunto completo dell'opportunità di investimento. Essendo una misura \gls{scale-free}, 
permette di confrontare in modo diretto strumenti finanziari differenti. 
Inoltre, le serie storiche dei rendimenti presentano proprietà statistiche favorevoli, prima tra tutte la \gls{stazionarieta}, che agevola 
l'identificazione di modelli predittivi e lo studio della volatilità.

Per misurare la variazione del prezzo nel tempo si distinguono due tipi di rendimento: il rendimento semplice uniperiodale e il 
rendimento composto a tempo continuo (log-rendimento).


\subsection{Rendimento semplice uniperiodale}
Sia $P_t$ il prezzo di un asset alla data $t$ e $P_{t-1}$ il suo prezzo nel periodo precedente, e si assuma che l'asset non paghi dividendi\cite{campbell:econometrics}.
Il rendimento netto semplice uniperiodale $R_t$, nell'intervallo di tempo $[t - 1, t]$, è definito come:
\[
    R_t = \frac{P_t}{P_{t-1}} - 1 = \frac{P_t - P_{t-1}}{P_{t-1}}
\]
e di conseguenza il rendimento lordo semplice, cioè il capitale totale dopo la detenzione dell'asset per un periodo è pari a $1 + R_t = \frac{P_t}{P_{t-1}}$.

Partendo da questo è possibile estendere la misurazione su orizzonti temporali più ampi, calcolando il rendimento lordo di un asset sugli ultimi $k$ periodi.
Denotato come $1 + R_t(k)$, equivale al prodotto dei $k$ rendimenti uniperiodali, ovvero:
\begin{equation*}
\begin{split}
    1 + R_t(k) &= (1 + R_t) \cdot (1 + R_{t-1}) \cdots (1 + R_{t-k+1}) \\
               &= \frac{P_t}{P_{t-1}} \cdot \frac{P_{t-1}}{P_{t-2}} \cdots \frac{P_{t-k+1}}{P_{t-k}} \\
               &= \frac{P_t}{P_{t-k}}
\end{split}
\end{equation*}
Il corrispondente rendimento netto multiperiodale si ottiene sottraendo l'unità: $R_t(k) = \frac{P_t}{P_{t-k}} - 1$.

Questo mostra che il rendimento semplice non gode della proprietà di additività temporale, rendendo necessaria la moltiplicazione. Per questo motivo si preferisce utilizzare i log-rendimenti.


\subsection{Rendimento composto a tempo continuo}
Un'alternativa al rendimento semplice uniperiodale è rappresentata dal rendimento composto a tempo continuo, noto anche come log-rendimento.

Il rendimento composto a tempo continuo calcolato su un singolo periodo $r_t$ è definito come il logaritmo naturale del rendimento lordo semplice:
\[
    r_t = \ln(1 + R_t) = \ln\left(\frac{P_t}{P_{t-1}}\right)
\]

Definendo $p_t = \ln(P_t)$, attraverso le proprietà dei logaritmi la formula si semplifica in: 
\[
    r_t = p_t - p_{t-1}
\]

Il vantaggio emerge nel calcolo su orizzonti temporali ampi, in quanto il log-rendimento multiperiodale $r_t(k)$ è calcolabile come somma dei $k$ log-rendimenti:
\begin{align*}
    r_t(k) &= \ln(1 + R_t(k)) \\
           &= \ln(1 + R_t) + \ln(1 + R_{t-1}) + \dots + \ln(1 + R_{t-k+1}) \\
           &= r_t + r_{t-1} + \dots + r_{t-k+1}
\end{align*}

I rendimenti composti a tempo continuo multiperiodali corrispondono alla somma dei rendimenti composti a tempo continuo uniperiodali e ciò li rende preferibili per l'analisi econometrica e lo studio della volatilità.


\section{La volatilità}
\label{sec:volatilita}
In ambito finanziario la volatilità non coincide con il rischio, ma se interpretata come incertezza diventa un elemento chiave per le decisioni 
di investimento e per la creazione dei portafogli\cite{poon:forecasting}. Infatti, il punto di partenza per la valutazione del rischio 
è una corretta previsione della futura volatilità dei prezzi degli asset.

Per via del forte legame tra l'incertezza dei mercati e la fiducia pubblica, le stime della volatilità vengono usate dai responsabili di 
politica economica come «barometro» per misurare la vulnerabilità dei mercati e le banche centrali la monitorano per prendere decisioni 
di politica monetaria. Per questo motivo la volatilità dei mercati finanziari può avere ampie ripercussioni sull'intera economia.

Una caratteristica problematica della volatilità azionaria è che essa, a differenza dei prezzi o dei rendimenti giornalieri, non è direttamente osservabile e rimane una \gls{variabile-latente}\cite{tsay:analysis}.
Non è infatti possibile estrapolare la volatilità a partire da un singolo rendimento giornaliero.
Per ovviare a questo problema, si ricorre spesso all'uso di misure \gls{proxy}, come il valore assoluto dei rendimenti ($|r_t|$) o i rendimenti al quadrato ($r_t^2$) per ottenere una misura della varianza.
Tuttavia, la necessità di ricorrere a tali approssimazioni, data l'impossibilità di osservare il dato reale, rende complessa la valutazione delle performance dei modelli previsionali.

Sebbene non sia direttamente osservabile, essa presenta alcune caratteristiche riscontrabili nei rendimenti degli asset, come la presenza di \emph{volatility clusters} (fasi di alta volatilità alternate a periodi di calma) e una reazione asimmetrica agli shock, nota come effetto leva.

Queste caratteristiche, definibili come «fatti stilizzati», evidenziano le debolezze dei modelli tradizionali e rendono necessaria l'adozione di modelli avanzati per catturare la volatilità condizionata.


\section[Fatti stilizzati]{Fatti stilizzati delle serie storiche finanziarie}
Le proprietà statistiche che emergono regolarmente nello studio delle serie storiche, comuni a diversi strumenti, mercati e orizzonti temporali 
prendono il nome di fatti stilizzati\cite{cont:empirical}.

Conoscere tali caratteristiche è necessario per costruire modelli econometrici adeguati. I modelli finanziari tradizionali, infatti, 
si basano sull'assunzione di normalità dei rendimenti e varianza costante nel tempo (\gls{omoschedasticita}). Tuttavia, l'osservazione dei fatti 
stilizzati invalida queste ipotesi di base, rendendo i modelli lineari classici inadeguati a misurare correttamente il rischio.

Per giustificare l'adozione di modelli avanzati in grado di catturare l'evoluzione della volatilità condizionata, nelle sezioni seguenti 
vengono illustrate le principali proprietà statistiche osservate nelle serie finanziarie, con particolare attenzione alla stazionarietà, 
alla leptocurtosi, al \emph{volatility clustering} e all'effetto leva.


\subsection{Stazionarietà}
La stazionarietà è un requisito di base per l'analisi delle serie storiche: un processo è stazionario quando media, varianza e struttura di 
autocorrelazione restano costanti nel tempo. Quando queste proprietà cambiano nel tempo, la serie è non stazionaria, spesso a causa della 
presenza di \emph{trend} o radici unitarie. Tale condizione si verifica empiricamente tramite appositi test, come quello di Dickey e 
Fuller\cite{dickey:estimators} e la sua estensione, l'\gls{adf}.

Per i rendimenti finanziari, la stazionarietà garantisce l'esistenza di una varianza incondizionata finita, attorno alla quale si muove la 
volatilità condizionata. A seguito di uno shock sui prezzi, l'incertezza non aumenta senza limite, ma tende a tornare progressivamente verso 
il suo livello medio di lungo periodo.

Nei modelli GARCH che verranno trattati, Bollerslev\cite{bollerslev:garch} mostra che, con coefficienti non negativi, il processo è stazionario 
in senso lato quando la somma dei coefficienti ARCH e GARCH è minore di 1. In questo caso la varianza incondizionata è costante e finita 
e gli shock di mercato hanno un effetto solo temporaneo sulla volatilità.


\subsection{Leptocurtosi}
La curtosi è un indicatore utile per descrivere la forma di una distribuzione di rendimenti: misura il peso delle code rispetto alla normale, 
cioè la frequenza con cui si manifestano valori estremi\cite{tsay:analysis}.

Mentre una distribuzione normale standard presenta una curtosi teorica pari a 3, le serie storiche finanziarie mostrano tipicamente una 
struttura leptocurtica. Ciò si traduce in un eccesso di curtosi positivo ($K - 3 > 0$): le code risultano più pesanti e i valori estremi, 
sia positivi che negativi, occorrono con una probabilità superiore a quella prevista dal modello gaussiano\cite{cont:empirical}.

In passato si tendeva a considerare le grandi variazioni come \emph{outlier} generati da meccanismi speciali, cercando modi per escluderle 
dal campione\cite{fama:behavior}.
Questi movimenti, però, non possono essere scartati, perché possono pesare molto sul rendimento complessivo di un periodo e vanno 
modellati con attenzione, ad esempio per calcolare il \gls{var}\cite{cont:empirical}.

La natura non gaussiana dei dati rende insufficiente l'uso della deviazione standard per misurare la dispersione\cite{cont:empirical}. 
Inoltre, assumere la normalità dei rendimenti porta a sottostimare la probabilità degli eventi estremi; per questo, prima di applicare 
modelli per la volatilità come gli ARCH\cite{engle:arch}, conviene verificare la non normalità dei rendimenti.

La presenza di code pesanti, suggerita da una curtosi maggiore di 3, viene poi accertata tramite test statistici, come il test di Jarque-Bera, 
che valuta congiuntamente asimmetria e curtosi\cite{jarque:bera}, e con strumenti grafici come il QQ-plot.

I modelli GARCH risultano adatti all'analisi finanziaria perché, come i modelli ARCH, generano una distribuzione incondizionata 
con code più pesanti di quelle della normale, anche quando gli errori sono gaussiani\cite{tsay:analysis}.

% immagine
\begin{figure}[H]
    \centering
    \includegraphics[width=0.8\textwidth]{QQ_Plot_S_P_500_SPY_.png}
    \caption[QQ Plot dei rendimenti dell'indice S\&P 500]{QQ Plot dei rendimenti dell'indice S\&P 500 (rispetto alla distribuzione normale).}
    \label{fig:qqplot-sp500}
\end{figure}


\subsection{Volatility clustering}
Il \emph{volatility clustering} (volatilità a grappoli) è un fenomeno che caratterizza le serie storiche finanziarie. 
Come osservato da Mandelbrot\cite{mandelbrot:variation}, le grandi variazioni di prezzo tendono ad essere seguite da altre grandi variazioni 
mentre piccole fluttuazioni da variazioni altrettanto contenute, indipendentemente dal segno, mostrando che eventi ad alta volatilità tendono 
a raggrupparsi nel tempo.

L'assenza di autocorrelazione nei rendimenti lineari non implica assoluta indipendenza, al contrario funzioni non lineari semplici, 
come rendimenti al quadrato o in valore assoluto, mostrano un'autocorrelazione positiva. 
Questa autocorrelazione positiva prolungata è coerente con la presenza di \emph{volatility clustering}\cite{cont:empirical}.

Come già osservato nel Paragrafo \ref{sec:volatilita} (\nameref{sec:volatilita}), la volatilità giornaliera non è direttamente osservabile dai soli rendimenti di chiusura, 
il che rende complesso lo studio di questa dinamica. 
I modelli tradizionali risultano inadatti a questo scopo poiché assumono varianza condizionata costante nel tempo.

Per superare questo limite Engle\cite{engle:arch} introdusse la classe dei processi ARCH, passando al concetto di varianza condizionata 
($\sigma_t^2$), ovvero una varianza che varia nel tempo sulla base delle informazioni passate. 
Come descritto in Tsay\cite{tsay:analysis}, l'estensione ai modelli GARCH permette di catturare questo fenomeno, dato che un ampio valore 
dei rendimenti passati al quadrato genera grande varianza nel periodo successivo, replicando così la natura a grappoli della volatilità 
osservata nei mercati.


\subsection{Asimmetria ed effetto leva}
Un'altra caratteristica ricorrente delle serie storiche finanziarie è l'asimmetria della volatilità in risposta agli shock. 
Come evidenziato da Cont\cite{cont:empirical}, i rendimenti sono correlati negativamente con la variazione della volatilità futura: 
questo fenomeno è noto come effetto leva (\emph{leverage effect}) ed è stato inizialmente documentato da Black\cite{black:studies}.

Come riassunto da Nelson\cite{nelson:conditional}, la volatilità tende ad aumentare dopo cattive notizie e a diminuire dopo buone notizie. 
Nonostante l'efficacia nel catturare il fenomeno di \emph{volatility clustering}, i modelli GARCH tradizionali risultano inadatti in quanto 
assumono che variazioni positive e negative della stessa entità producano lo stesso impatto sulla varianza futura\cite{tsay:analysis}.

Per superare questo problema di simmetria e catturare l'effetto leva sono state sviluppate estensioni asimmetriche, come il GJR-GARCH e 
l'EGARCH, che vengono trattate nel capitolo successivo. %2
    % capitolo 2: Modelli a volatilità condizionata
\chapter{Modelli a volatilità condizionata}
\label{cap:volatilita-condizionata}

 %3
    % capitolo 3: Dati e metodologia
\chapter{Dati e metodologia}
\label{cap:dati-metodologia}


% \section{Prerequisiti}

 %4
    % capitolo 5: Risultati dei modelli GARCH e delle correlazioni
\chapter{Risultati dei modelli GARCH e delle correlazioni}
\label{cap:risultati}

\section{Selezione del modello e stima dei parametri}
L'analisi empirica ha inizio con la selezione del modello ottimale per analizzare le dinamiche della varianza condizionata, per poi passare alla stima e interpretazione 
dei parametri nel modello selezionato. La valutazione si basa sulla minimizzazione del Criterio di Informazione Bayesiano (BIC), confrontando le strutture sGARCH, 
GJR-GARCH ed EGARCH, stimate con distribuzione Gaussiana e Student-t. Come mostrato nella tabella \ref{tab:bic_valori}, il modello asimmetrico EGARCH-t minimizza il valore 
per tutti e 6 gli asset in esame.

\begin{table}[htbp]
    \centering
    \caption{Valori del BIC}
    \label{tab:bic_valori}
    \resizebox{\textwidth}{!}{%
    \begin{tabular}{lcccccc}
        \toprule
        \textbf{Asset} & \textbf{sGARCH\_norm} & \textbf{GJR-GARCH\_norm} & \textbf{EGARCH\_norm} & \textbf{sGARCH\_std} & \textbf{GJR-GARCH\_std} & \textbf{EGARCH\_std} \\
        \midrule
        Oro & 2.6817 & 2.6722 & 2.6677 & 2.6190 & 2.6152 & 2.6086 \\
        Nucor & 4.3930 & 4.3903 & 4.3759 & 4.3039 & 4.3041 & 4.2960 \\
        GM & 4.3607 & 4.3561 & 4.3307 & 4.2492 & 4.2453 & 4.2368 \\
        Apple & 3.8663 & 3.8441 & 3.8228 & 3.7610 & 3.7500 & 3.7397 \\
        S\&P 500 (SPY) & 2.6115 & 2.5867 & 2.5623 & 2.5374 & 2.5121 & 2.5007 \\
        USD/CNH & 0.2254 & 0.2214 & 0.2113 & 0.1146 & 0.1157 & 0.1091 \\
        \bottomrule
    \end{tabular}%
    }
\end{table}

Questa superiorità viene confermata anche dal Criterio di Informazione di Akaike (AIC) e dalle tabelle delta (dBIC e dAIC), i cui risultati completi sono riportati in 
Appendice. Impostando l'EGARCH-t come benchmark ($\Delta=0$), tutti gli altri modelli restituiscono delta ampiamente positivi e mostrano come la Student-t migliori 
l'adeguatezza in tutti i modelli. Definita la specifica ottimale, la Tabella \ref{tab:parametri_egarch} riporta le stime dei parametri del modello EGARCH-t.

\begin{table}[htbp]
    \centering
    \caption{Stima dei parametri del modello EGARCH-t}
    \label{tab:parametri_egarch}
    \resizebox{\textwidth}{!}{%
    \begin{tabular}{lcccccc}
    \toprule
    \textbf{Parametro} & \textbf{Oro} & \textbf{Nucor} & \textbf{GM} & \textbf{Apple} & \textbf{S\&P 500 (SPY)} & \textbf{USD/CNH} \\
    \midrule
    $\mu$ & 0.059*** & 0.068* & 0.058** & 0.112*** & 0.070*** & -0.001 \\
    $\omega$ & -0.002 & 0.014*** & 0.013*** & 0.031*** & -0.018*** & -0.057*** \\
    $\theta^*$ (segno) & 0.036*** & -0.023** & -0.043*** & -0.092*** & -0.165*** & 0.019 \\
    $\alpha^*$ (ampiezza) & 0.102*** & 0.088*** & 0.101*** & 0.175*** & 0.191*** & 0.200*** \\
    $\beta^*$ & 0.988*** & 0.990*** & 0.992*** & 0.968*** & 0.970*** & 0.979*** \\
    $\nu$ & 5.353*** & 5.261*** & 4.157*** & 4.622*** & 6.099*** & 4.445*** \\
    \bottomrule
    \end{tabular}%
    }
    \vspace{0.15cm}
    \footnotesize{\textit{Nota:} *, **, *** indicano significatività al 10\%, 5\% e 1\% (basata su errori standard classici). Le stime si riferiscono al modello includente le variabili dummy, i cui coefficienti sono omessi per brevità; il $\beta^*$ riportato è quello impiegato per i calcoli nella Sezione 5.2. Rispetto a \textit{rugarch}: $\theta^*$ = alpha1, $\alpha^*$ = gamma1, $\omega$ = costante.}
\end{table}

Il parametro di asimmetria ($\theta^*$) evidenzia il classico effetto leva per i titoli azionari (GM e Apple) e per l'indice di mercato S&P 500, mentre per Nucor risulta 
significativo solo al 10\%. I valori inferiori a 0 indicano che gli shock negativi incrementano la volatilità in misura maggiore rispetto a shock positivi di pari entità. 
Per quantificare la dimensione di questa asimmetria, si osservi l'indice S&P 500: l'impatto dei rialzi pesa solamente 0,026 ($\alpha^* + \theta^*$), mentre quello dei ribassi 
pesa ben 0,356 ($\alpha^* - \theta^*$). Al contrario, l'Oro presenta un coefficiente positivo (0,036), evidenziando un'asimmetria inversa, diversa da quella azionaria. 
La valutazione della sua effettiva natura di bene rifugio viene pertanto demandata all'analisi delle correlazioni nei periodi di stress, mentre per il tasso di cambio USD/CNH 
non si rifiuta la simmetria. 

I gradi di libertà della distribuzione ($\nu$) si attestano tra 4 e 6, confermando la presenza di code grasse (\emph{fat tails}) e l'inadeguatezza di una distribuzione 
normale. Infine, il coefficiente di persistenza ($\beta^*$) risulta prossimo a 1 per tutti gli asset (superiore a 0.96). Calcolando l'\gls{emivita} (\emph{half-life}) degli 
shock tramite la formula $\log(0.5)/\log(\beta^*)$, emerge che i movimenti si dimezzano in tempi variabili: da un assorbimento più rapido di circa 21 giorni per Apple, 
a un rientro lento di quasi 86 giorni per General Motors, evidenziando un'elevata persistenza della volatilità. %5
    % capitolo 6: conclusioni
\chapter{Conclusioni}
\label{cap:conclusioni}

L'elaborato ha valutato se gli annunci tariffari tra Stati Uniti e Cina del biennio 2025-2026 hanno modificato la volatilità condizionata dei 
sei strumenti finanziari selezionati, con quale intensità, persistenza e con quali differenze tra asset class. L'analisi si è basata su un modello 
EGARCH(1,1) con errori t di Student, esteso con variabili \emph{dummy} definite su una finestra di due sedute per ogni evento selezionato.

Il modello minimizza il BIC per tutti e sei gli asset rispetto alle altre specificazioni considerate (sGARCH, GJR-GARCH ed EGARCH, con errori gaussiani e t di Student). 
I parametri stimati confermano i fatti stilizzati discussi nel Capitolo 2: code pesanti (da 4,2 a 6,1 gradi di libertà), persistenza elevata della volatilità ($\beta^* > 0,96$) 
e un effetto leva significativo per S\&P 500, Apple, GM e Nucor. Per l'oro l'asimmetria è invece di segno opposto, mentre per l'USD/CNH non è significativa. Poiché l'emivita degli shock va da 
circa 21 sedute per Apple a circa 86 per General Motors, l'effetto di un annuncio si prolunga per settimane oltre le due sedute della finestra.

Il test congiunto fornisce un'evidenza differenziata: per l'oro l'ipotesi nulla è rifiutata anche dopo la correzione di Holm, per Apple 
l'evidenza è intermedia (significativa al 10\%), per S\&P 500, USD/CNH e Nucor il rifiuto al 5\% vale solo prima della correzione, 
mentre per GM non emerge alcuna significatività. Il BIC preferisce sempre il modello senza \emph{dummy}, ma ciò dipende dalla penalizzazione 
dei sette parametri aggiuntivi e non è stato usato per escluderle.

Il risultato più solido riguarda il \emph{Liberation Day} (2 aprile 2025), che aumenta in modo significativo la volatilità di cinque asset 
su sei, mentre per GM l'effetto è positivo ma non significativo. Un secondo episodio rilevante è la crisi delle terre rare, con aumenti 
significativi per oro, Nucor e S\&P 500. Gli eventi di distensione si associano a riduzioni meno frequentemente significative, mentre per 
il vertice di Ginevra, il provvedimento anti-Fentanyl e la sentenza della Corte Suprema non emergono effetti stabili. I coefficienti negativi 
dell'escalation del 9 aprile non indicano un calo dell'incertezza: il modello prevede già una volatilità molto elevata dopo il \emph{Liberation Day}, e il 
coefficiente ne corregge l'effetto. Gli annunci di escalation più rilevanti (Liberation Day e terre rare) sembrano associarsi ad aumenti di volatilità più ampi e 
frequenti di quelli di distensione, e la reazione agli annunci tra gli strumenti è eterogenea: GM non mostra effetti significativi congiunti nonostante l'esposizione 
diretta agli annunci, mentre l'oro, l'S\&P 500 e Nucor reagiscono alla crisi delle terre rare in modo diverso da Apple e dal cambio. La diagnostica sui 
residui standardizzati non evidenzia criticità rilevanti.

L'analisi delle correlazioni tra i residui standardizzati mostra che i legami tra le coppie di asset analizzate sono cambiati 
in modo eterogeneo nel periodo dei dazi. La correlazione tra S\&P 500 e oro sale da 0,028 a 0,166, e la correlazione \emph{rolling} raggiunge 
il massimo dell'intero campione nel 2026: in una fase del periodo di dazi oro e azionario sembrano aver risposto allo stesso modo a shock comuni, 
con una possibile riduzione del beneficio di diversificazione dell'oro nei momenti di stress. Per Nucor e GM la correlazione scende da 0,396 a 0,220 e la 
curva \emph{rolling} si avvicina a zero, coerentemente con un impatto differenziato delle tariffe sui due settori, che secondo l'ipotesi discussa nella Sezione \ref{sec:definizione_asset} 
favorirebbero i produttori siderurgici statunitensi e penalizzerebbero GM. Per S\&P 500 e USD/CNH la correlazione media è quasi invariata, ma questo nasconde una forte 
variazione interna, in quanto la correlazione \emph{rolling} scende fino a circa $-0,7$ nel 2026, a indicare che in quella fase i ribassi dell'azionario statunitense coincidono più 
spesso con un indebolimento dello yuan.

In conclusione, lo studio mostra come gli annunci tariffari tra Stati Uniti e Cina siano associati a variazioni della volatilità 
dei mercati finanziari, in misura eterogenea tra strumenti ed eventi, e sottolinea l'importanza di modelli in grado di cogliere la 
volatilità condizionata. I risultati vanno letti come evidenza di un'associazione, ma non implicano una diretta relazione di causa-effetto: 
ogni evento è identificato da due sole osservazioni e i test sono approssimati. Come possibili sviluppi futuri, si considerano l'estensione 
dell'orizzonte temporale, l'uso di finestre di evento di ampiezza diversa e il confronto con ulteriori specificazioni, oltre all'inclusione 
di ulteriori strumenti, mercati ed eventi di natura economica e geopolitica.

%\section{Consuntivo finale}

%\section{Raggiungimento degli obiettivi}

%\section{Conoscenze acquisite}

%\section{Valutazione personale}
 %6

    \appendix
    \cleardoublepage
\phantomsection
\addcontentsline{toc}{chapter}{Appendice}
\chapter*{Appendice}
\markboth{Appendice}{}
\setcounter{chapter}{1} % le eventuali sezioni si chiameranno A.1, A.2, ...
\setcounter{figure}{0} % le figure si chiamano A.1, A.2, ...

% QQ Plot in riferimento al capitolo 2 (Analisi delle serie storiche finanziarie)
\section{QQ plot dei rendimenti}
\label{app:qqplot}
Di seguito sono riportati i QQ plot dei rendimenti degli strumenti analizzati che non sono stati inseriti nel Capitolo~\ref{cap:analisi-serie}, 
dove è presentato quello dell'indice S\&P 500 (Figura~\ref{fig:qqplot-sp500}).

\begin{figure}[H]
    \centering
    \includegraphics[width=0.75\textwidth]{QQ_Plot_Oro.png}
    \caption[QQ Plot dei rendimenti dell'oro]{QQ Plot dei rendimenti dell'oro rispetto alla distribuzione normale.}
    \label{fig:qqplot-oro}
\end{figure}

\begin{figure}[H]
    \centering
    \includegraphics[width=0.75\textwidth]{QQ_Plot_USD_CNH.png}
    \caption[QQ Plot dei rendimenti del tasso di cambio USD/CNH]{QQ Plot dei rendimenti del tasso di cambio USD/CNH rispetto alla distribuzione normale.}
    \label{fig:qqplot-usdcnh}
\end{figure}

\begin{figure}[H]
    \centering
    \includegraphics[width=0.75\textwidth]{QQ_Plot_Nucor.png}
    \caption[QQ Plot dei rendimenti di Nucor]{QQ Plot dei rendimenti di Nucor rispetto alla distribuzione normale.}
    \label{fig:qqplot-nucor}
\end{figure}

\begin{figure}[H]
    \centering
    \includegraphics[width=0.75\textwidth]{QQ_Plot_GM.png}
    \caption[QQ Plot dei rendimenti di General Motors]{QQ Plot dei rendimenti di General Motors rispetto alla distribuzione normale.}
    \label{fig:qqplot-gm}
\end{figure}

\begin{figure}[H]
    \centering
    \includegraphics[width=0.75\textwidth]{QQ_Plot_Apple.png}
    \caption[QQ Plot dei rendimenti di Apple]{QQ Plot dei rendimenti di Apple rispetto alla distribuzione normale.}
    \label{fig:qqplot-apple}
\end{figure}

\newpage
% immagini in riferimento al capitolo 4 (Dati e metodologia)
\section{Prezzi di chiusura e log-rendimenti}
\label{app:prezzi}
Di seguito sono riportati i prezzi di chiusura giornalieri e i log-rendimenti che non sono stati inseriti nel capitolo~\ref{cap:dati-metodologia}, 
dove sono stati mostrati solo S\&P 500, Apple e USD/CNH.

\begin{figure}[H]
    \centering
    \begin{subfigure}{0.48\textwidth}
        \includegraphics[width=\linewidth]{Prezzo_Oro.png}
    \end{subfigure}\hfill
    \begin{subfigure}{0.48\textwidth}
        \includegraphics[width=\linewidth]{Rendimenti_Oro.png}
    \end{subfigure}

    \medskip
    \begin{subfigure}{0.48\textwidth}
        \includegraphics[width=\linewidth]{Prezzo_Nucor.png}
    \end{subfigure}\hfill
    \begin{subfigure}{0.48\textwidth}
        \includegraphics[width=\linewidth]{Rendimenti_Nucor.png}
    \end{subfigure}

    \medskip
    \begin{subfigure}{0.48\textwidth}
        \includegraphics[width=\linewidth]{Prezzo_GM.png}
    \end{subfigure}\hfill
    \begin{subfigure}{0.48\textwidth}
        \includegraphics[width=\linewidth]{Rendimenti_GM.png}
    \end{subfigure}
    \caption[Prezzi di chiusura e log-rendimenti]{Prezzi di chiusura giornalieri (a sinistra) e log-rendimenti giornalieri (a destra) per oro, Nucor e General Motors.}
    \label{fig:app-prezzi-rendimenti}
\end{figure}
 

\newpage
\section{Funzione di autocorrelazione (ACF) e ACF dei residui al quadrato}
\label{app:acf}
Di seguito sono riportate le ACF dei log-rendimenti e le ACF dei residui al quadrato ($r_t^2$) degli asset analizzati, a integrazione della Figura~\ref{fig:acf-plot}.

\begin{figure}[H]
    \centering
    \begin{subfigure}{0.48\textwidth}
        \includegraphics[width=\linewidth]{ACF_Oro.png}
    \end{subfigure}\hfill
    \begin{subfigure}{0.48\textwidth}
        \includegraphics[width=\linewidth]{ACF_Sq_Oro.png}
    \end{subfigure}

    \medskip

    \begin{subfigure}{0.48\textwidth}
        \includegraphics[width=\linewidth]{ACF_Nucor.png}
    \end{subfigure}\hfill
    \begin{subfigure}{0.48\textwidth}
        \includegraphics[width=\linewidth]{ACF_Sq_Nucor.png}
    \end{subfigure}

    \caption[ACF dei log-rendimenti e ACF dei residui al quadrato (Oro e Nucor)]{Funzione di Autocorrelazione (ACF) dei log-rendimenti e dei residui al quadrato ($r_t^2$) per Oro e Nucor. Le linee tratteggiate indicano le bande di confidenza al 95\%.}
    \label{fig:app-acf-1}
\end{figure}

\begin{figure}[H]
    \centering
    \begin{subfigure}{0.48\textwidth}
        \includegraphics[width=\linewidth]{ACF_GM.png}
    \end{subfigure}\hfill
    \begin{subfigure}{0.48\textwidth}
        \includegraphics[width=\linewidth]{ACF_Sq_GM.png}
    \end{subfigure}

    \medskip

    \begin{subfigure}{0.48\textwidth}
        \includegraphics[width=\linewidth]{ACF_Apple.png}
    \end{subfigure}\hfill
    \begin{subfigure}{0.48\textwidth}
        \includegraphics[width=\linewidth]{ACF_Sq_Apple.png}
    \end{subfigure}

    \caption[ACF dei log-rendimenti e ACF dei residui al quadrato (GM e Apple)]{Funzione di Autocorrelazione (ACF) dei log-rendimenti e dei residui al quadrato ($r_t^2$) per GM e Apple. Le linee tratteggiate indicano le bande di confidenza al 95\%.}
    \label{fig:app-acf-2}
\end{figure}

\newpage
% altre immagini non direttamente presenti nel capitolo 4, in cui vengono mostrati tutti gli asset analizzati
\section{Funzione di autocorrelazione parziale (PACF)}
\label{app:pacf}
Di seguito sono riportate le funzioni di autocorrelazione parziale (PACF) dei log-rendimenti di tutti gli asset analizzati.
 
\begin{figure}[H]
    \centering
    \includegraphics[width=\textwidth]{PACF_Tutti_Appendice.png}
    \caption[PACF dei log-rendimenti]{Funzione di Autocorrelazione Parziale (PACF) dei log-rendimenti. Le linee tratteggiate indicano le bande di confidenza al 95\%.}
    \label{fig:app-pacf}
\end{figure}


\newpage
\section{Proxy della varianza}
\label{app:proxy}
Di seguito è riportato l'andamento dei quadrati dei log-rendimenti ($r_t^2$), utilizzati come proxy della varianza, per tutti gli asset analizzati.
 
\begin{figure}[H]
    \centering
    \includegraphics[width=\textwidth]{Proxy_Varianza_Tutti_Appendice.png}
    \caption[Proxy della varianza]{Quadrati dei log-rendimenti ($r_t^2$) come proxy della varianza.}
    \label{fig:app-proxy}
\end{figure}


\newpage
% immagini in riferimento al capitolo 5
\section{Volatilità condizionata con e senza \emph{dummy}}
\label{app:sigma}
Di seguito è riportata la volatilità condizionata stimata con e senza le variabili \emph{dummy} per tutti gli asset analizzati.
 
\begin{figure}[H]
    \centering
    \includegraphics[width=\textwidth]{fig_sigmaApp.png}
    \caption[Volatilità condizionata con e senza dummy]{Volatilità condizionata ($\sigma_t$, in percentuale e scala logaritmica) di Nucor, GM e USD/CNH, 
    stimata con e senza le variabili \emph{dummy}. Le linee verticali numerate indicano gli eventi tariffari, nello stesso ordine della Sezione~\ref{sec:eventi}.}
    \label{fig:app-sigma}
\end{figure}

    \backmatter
    % \glsaddall % per stampare tutto anche se i termini non ci sono ancora nel testo

    % fix intestazione acronimi e glossario
    \cleardoublepage
    % IMPONE IL NOME ACRONIMI o GLOSSARIO FISSO A SINISTRA NELL'INTESTAZIONE E NULLA A DESTRA
    \fancyhead[L]{\fontsize{12}{14.5} \selectfont Acronimi e abbreviazioni}
    \fancyhead[R]{} 
    \printglossary[type=\acronymtype, title=Acronimi e abbreviazioni, toctitle=Acronimi e abbreviazioni]

    \cleardoublepage
    \fancyhead[L]{\fontsize{12}{14.5} \selectfont Glossario}
    \printglossary[type=main, title=Glossario, toctitle=Glossario]

    % STESSA COSA PER BIBLIOGRAFIA E RINGRAZIAMENTI (se dovessero occupare più di una pagina)
    \cleardoublepage
    \fancyhead[L]{\fontsize{12}{14.5} \selectfont Bibliografia}
    \cleardoublepage
\chapter{Bibliografia}

\nocite{*}

\printbibliography[heading=subbibliography,title={Libri, testi e articoli},category=opere]
\printbibliography[heading=subbibliography,title={Sitografia},category=web]

    %\cleardoublepage
\clearpage
\phantomsection
\pdfbookmark{Ringraziamenti}{ringraziamenti}

\pagestyle{empty}

\chapter*{Ringraziamenti}
\thispagestyle{empty}
\begingroup
\setlength{\parskip}{0.6\baselineskip}
\setlength{\parindent}{0pt}

% testo ringraziamenti




\bigskip
\noindent\textit{\myLocation, \myTime}
\hfill \myName

\cleardoublepage
\pagestyle{fancy} % ringraziamenti
\end{document}